% Section 4.6, from lectures/L04/appendix/e_exercises.md.

\section{Exercises}
\label{c4:sec:exercises}

\begin{aside}
\textbf{How to check your work.} Exercise~\ref{c4:ex:array} is checked by this chapter's test
suite: write the file at the path the specification gives, then run \code{make test}. The suite is
cumulative and runs Chapter~1 to Chapter~3's tests too.

The rest will have worked solutions in the companion repository, published later as
\file{lectures/L04/appendix/f_solutions.md}, and the hand calculations have short answers in
Appendix~\ref{app:answers}. Work each one before you read it.

\textbf{No hardware is needed for any of this.}
\end{aside}

Each exercise is labelled with its kind. There is exactly one \textbf{Cross-check}, and here it is
Exercise~\ref{c4:ex:crosscheck}.

% ------------------------------------------------------------------------------------------
\exercise{Three pointers}{Recall}
\label{c4:ex:pointers}
\begin{parts}
  \item Name the three pointer register pairs and the registers each is made of, and say which
    register of a pair holds the high byte.
  \item Which of the three has no displacement form, and which one does compiled C reserve?
  \item \code{movw r30, r24} copies a pointer in one instruction. Say why one instruction is enough,
    and what would be needed if the two pairs did not agree about byte order.
  \item Nothing stops you putting a loop counter in \code{r30}. Give the situation in which that
    turns out to be expensive.
\end{parts}

% ------------------------------------------------------------------------------------------
\exercise{Four ways to read a byte}{Hand calculation}
\label{c4:ex:fourways}
\code{Z} holds \code{0x0204}. The bytes at \code{0x0202} to \code{0x0208} are, in order:
\code{11 22 33 44 55 66 77}.

For each of the following, give the value loaded into \code{r16} and the value of \code{Z}
afterwards:
\begin{parts}
  \item \code{ld r16, Z}
  \item \code{ld r16, Z+}
  \item \code{ld r16, -Z}
  \item \code{ldd r16, Z+3}
  \item Two of the four leave \code{Z} unchanged. Which, and why is that the property that makes
    structures cheap?
  \item All four cost the same. How many cycles, and what does that say about choosing between
    them?
\end{parts}

% ------------------------------------------------------------------------------------------
\exercise{Which pointer to use}{Design}
\label{c4:ex:which}
\begin{parts}
  \item A subroutine walks a buffer forwards, one byte at a time, and also has to read fields of a
    structure at known offsets. How many pointer registers does it need, which would you choose,
    and why not the third?
  \item A colleague's subroutine uses \code{Y} and does not save it. Say what breaks, when, and why
    the symptom appears in their caller rather than in their subroutine.
  \item Write down, in words, the four instructions a subroutine using \code{Y} has to add, and
    give their total cost in cycles.
  \item \code{ldd} cannot reach a displacement of 70. Give two ways to read that byte anyway, and
    say what each costs.
\end{parts}

% ------------------------------------------------------------------------------------------
\exercise{Array arithmetic}{Hand calculation}
\label{c4:ex:arrays}
An array of LED structures begins at \code{0x0300}.
\begin{parts}
  \item Give the address of entries 0, 1, 2 and 5.
  \item What is the stride, and where is it declared?
  \item A colleague advances the pointer by 6 instead. Give the address their code computes for
    entry 2, say which byte of which real entry that is, and describe what the driver would then
    read as a port register address.
  \item Their code produces no error and their first LED works. Explain both.
  \item An array of \emph{button} structures begins at the same address. Where is entry 3, and why
    is the answer different?
  \item How many LED structures fit in the ATmega328P's SRAM if nothing else used any of it? How
    many buttons? Say why neither number is the real limit on how many drivers a program can have.
\end{parts}

% ------------------------------------------------------------------------------------------
\exercise{The LED array}{Code}
\label{c4:ex:array}
Write \file{drivers/source/led_array.asm} to the specification in \secref{c4:sec:build}.
\begin{parts}
  \item Write all four subroutines.
  \item Run \code{make test}.
  \item Change your stride from \code{LED_SIZE} to 6, rebuild, and note which tests fail and which
    still pass. Then put it back.
  \item Move the count test in \code{led_array_init} from the top of the loop to the bottom,
    rebuild, and note which single test catches it. Explain why no other test could.
\end{parts}

% ------------------------------------------------------------------------------------------
\exercise{Where to put things}{Design}
\label{c4:ex:where}
\begin{parts}
  \item For each of \code{0x0005}, \code{0x0025}, \code{0x0068}, \code{0x0200}, \code{0x08FF} and
    \code{0x0900}, say what is there and whether a driver structure may live at it.
  \item A widely copied piece of AVR code allocates its structures at \code{RAMEND + 1}. Say what
    that address is on an ATmega328P, what a store to it does, and what the driver appears to do as
    a result.
  \item The same code on a different AVR, one with external RAM fitted, works. Explain.
  \item Give two ways to choose an address for a structure, and one advantage of each.
  \item Your program has one LED structure, one button structure, and an array of six LEDs. How
    many bytes of SRAM is that? What fraction of the total?
\end{parts}

% ------------------------------------------------------------------------------------------
\exercise{How deep does it go}{Hand calculation}
\label{c4:ex:depth}
Use the rule in \secref{c4:sec:worst}.
\begin{parts}
  \item A program with no interrupts, whose deepest path is a main loop calling a driver which
    calls \code{shift_bits}. How many bytes of stack?
  \item Chapter~3's program: a main loop that calls nothing, and an interrupt whose handler saves
    eight registers and \code{SREG} and calls \code{btn_pressed}, which itself calls
    \code{shift_bits}. How many bytes?
  \item Where does the stack pointer end up, and what is the lowest byte actually occupied? These
    are not the same number.
  \item A handler that is nothing but \code{reti} still costs stack. How much, and why?
  \item You add the LED array from Exercise~\ref{c4:ex:array}, and the main loop now calls
    \code{led_array_all_on}, which calls \code{led_on}, which calls \code{shift_bits}. Redo (b).
  \item A structure is placed at \code{0x08F0}. Using your answer to (b), is it safe? Show the
    comparison.
\end{parts}

% ------------------------------------------------------------------------------------------
\exercise[fillaccent]{How deep the stack really goes}{Cross-check}
\label{c4:ex:crosscheck}
\emph{Compute it by hand, measure it, reconcile.}
\begin{parts}
  \item \textbf{By hand.} Your answer to Exercise~\ref{c4:ex:depth}(b).
  \item \textbf{The lowest byte.} Turn that depth into an address: \code{SP} points at the next free
    byte, so the lowest byte the stack occupies is \code{RAMEND - depth + 1}. Do this before (c).
  \item \textbf{By measurement.} Let the program run, press the button, and read the low-water mark:
\begin{shell}
make measure IMAGE=app \
    CALLS="--run 4000 --drive B4=0 --run 300 --drive B4=1 --run 300"
\end{shell}
    \code{--run} lets the program run on its own, which is the only way to reach a handler: an
    interrupt cannot be called, it has to arrive. The last line reports how far down the stack
    pointer went.
  \item \textbf{Reconcile.} The byte count should agree with (a). The stack pointer's own value will
    be one below the lowest byte you computed in (b), and that is not a disagreement; say why.
  \item \textbf{Now run it without pressing anything.}
\begin{shell}
make measure IMAGE=app CALLS="--run 4000"
\end{shell}
    The measured depth is much smaller. Say what number you get, why it is smaller, and then answer
    the question this exercise is really about: \textbf{which of your three answers is now wrong,
    and which would you ship?}
  \item \textbf{The rule.} Write down, in one sentence, when a measurement of stack depth is worth
    more than the arithmetic and when it is worth less. Then say what a measured depth
    \emph{larger} than your computed one would mean, and which of the two you would go and check
    first.
  \item \textbf{The margin.} Take your worst-case depth and your structures from
    Exercise~\ref{c4:ex:where}(e). How many bytes are free? If a later change added one more level
    of nesting inside the handler, how much would that cost, and would you notice?
\end{parts}

% ------------------------------------------------------------------------------------------
\exercise{When it goes wrong}{Design}
\label{c4:ex:wrong}
\begin{parts}
  \item A wrong stride and an off-by-one loop bound are the two characteristic array bugs. For
    each, describe what memory looks like afterwards.
  \item Neither produces an error. Say what you would actually observe in a program with six LEDs
    and a wrong stride.
  \item The test suite catches both. Say what a test has to do that a program does not, in order to
    notice.
  \item A structure that ends at \code{0x08F0} and a stack that reaches 15 bytes down. Nothing fails for a
    week, and then it does. Give a plausible account of what changed.
  \item Name one thing you could add to the \emph{program} that would make a stack overflow
    noticeable rather than silent, and say what it would cost.
\end{parts}
